\subsection*{\texorpdfstring{$\lambda$-Calculus: Combinators}{-Calculus: Combinators}}

\textbf{\textit{Combinators}} are \ul{closed \iMbox{\lambda}-abstraction
    terms} which can be \ul{composed} to recreate \textit{other}
\ul{\iMbox{\lambda}-calculus} terms. A set of
\textbf{\textit{combinators}} which \ul{compose} to produce \ul{every}
other \ul{\iMbox{\lambda}-calculus} term \textit{(upto
    \ul{\iMbox{\beta \eta}-equivalence})} is called a
\textbf{\textit{complete combinatory basis}}. E.g. \iMbox{S,K,I} combinators.

\hSep
\vspace{2pt}

\paragraph*{Common Combinators:}

\kwub{Identity} => \iMbox{\mathrm{I} \equiv \lambda x.x}

\kwub{Identity once-removed} =>
\iMbox{\mathrm{I^{*}} \equiv \lambda fx.fx}, i.e. \stb{application}

\kwub{Thrush} => \iMbox{\mathrm{T} \equiv \lambda xf.fx},
i.e. \stb{reversed application}

\kwub{Kestrel} => \iMbox{\mathrm{K} \equiv \lambda xy.x}, i.e. \stb{constant}

\kwub{Kite} => \iMbox{\mathrm{KI} \equiv \lambda xy.y}

\kwub{Bluebird} => \iMbox{\mathrm{B} \equiv \lambda fgx.f(gx)}, i.e. \stb{composition}

\kwub{Queer bird} => \iMbox{\mathrm{Q \ or \ B'} \equiv \lambda fgx.g(fx)},
i.e. \stb{reversed composition}

\kwub{Mockingbird} => \iMbox{\omega  \equiv \lambda x. xx}
i.e. applying a function to itself

\kwub{Omega} => \iMbox{\Omega \equiv \omega \omega} is a \stu{non-terminating}
combinator which has no \stu{head normal form} \textit{(let alone
    \stu{$\beta$-normal form})} i.e. its \stb{unsolvable}

\kwub{Warbler} => \iMbox{\mathrm{W} \equiv \lambda fx. fxx}
i.e. \stb{argument duplicator}; \\
e.g. \iMbox{W \ (+) \ x \ \ \equiv \ \ x+x \ \ \equiv \ \  2x}

\kwub{Cardinal} => \iMbox{\mathrm{C} \equiv \lambda fxy. fyx}, i.e. \stb{flip}

\kwub{Viero} => \iMbox{\mathrm{V} \equiv \lambda xyz. zxy}, i.e. \stb{pair constructor}

\kwub{Starling} => \iMbox{\mathrm{S} \equiv \lambda fgx. (fx)(gx)} provides
\stu{context of \iMbox{x}} to \iMbox{f,g}; \textit{think: reader monad}

\kwub{Psi} => \iMbox{\Psi \equiv \lambda fgxy. f(gx)(gy)} processes
two arguments before combining them; \\
e.g. \iMbox{(\Psi \ (+) \ \text{length})} will add lengths of two lists

\hSep
\vspace{2pt}

\paragraph*{Recursion |> Fixpoint Combinators}

\kwub{fixed-Point} of a function is where it returns the unchanged input \\
i.e. for any \iMbox{x}, if \iMbox{f \ x = x} then \iMbox{x} is a \stu{fixed-point}
e.g. \iMbox{\mathrm{double} \ 0 = 0}

\kwub{fixed-Point combinator} \iMbox{\mathrm{fix}} is such that:
\begin{enumerate}
    \vItem If \iMbox{f} has \textit{any} \textbf{\textit{fixed-points}},
    then \iMbox{\mathrm{fix} \ f} is \textit{one} of those
    \textbf{\textit{fixed-points}},
    i.e. \iMbox{f \ (\mathrm{fix} \ f) =_{\beta} \mathrm{fix} \ f}

    \vItem These fixed point operators can be used to implement
    \textbf{\textit{recursion}} in \ul{$\lambda$-calculus},
    e.g. \ul{factorial} function \iMbox{\mathrm{fact}}:
    \begin{enumerate}
        \vItem
        \iMbox{\displaystyle \mathrm{fact} \equiv \lambda n. \mathrm{if} \ (\mathrm{isZero} \ n) \ \mathrm{then} \ 1 \ \mathrm{else} \ n \times (\mathrm{fact} \ (n-1) )}
        \vItem
        Now we can simply \textit{extract}
        \iMbox{\displaystyle \mathrm{fact} \equiv [\lambda f n. \mathrm{if} \ (\mathrm{isZero} \ n) \ \mathrm{then} \ 1 \ \mathrm{else} \ n \times (f \ (n-1) )] \ \mathrm{fact}}

        \begin{itemize}

            \vItem
            The inner part we can call \iMbox{F},
            i.e.~\iMbox{\displaystyle F \equiv \lambda f n. \mathrm{if} \ (\mathrm{isZero} \ n) \ \mathrm{then} \ 1 \ \mathrm{else} \ n \times (f \ (n-1) )}
            \vItem
            So now its
            \iMbox{\displaystyle \mathrm{fact} \equiv F \ \mathrm{fact}}
        \end{itemize}
        \vItem
        Now we use \iMbox{\mathrm{fix}} to re-write it as
        \iMbox{\mathrm{fact} \equiv \mathrm{fix} \ F}

        \begin{itemize}

            \vItem
            Since
            \iMbox{\displaystyle \mathrm{fix} \ F =_{\beta} F \ (\mathrm{fix} \ F)}
            and \iMbox{\mathrm{fix} \ F \equiv \mathrm{fact}}, we get
            \iMbox{\mathrm{fact} =_{\beta} F \ \mathrm{fact}}, i.e.~our
            \ul{original} definition
        \end{itemize}
    \end{enumerate}
\end{enumerate}


\hSep

\kwub{$Y$-combinator} => \iMbox{\mathrm{Y} 
\equiv \mathrm{B} \ \omega \ (\mathrm{C} \ \mathrm{B} \ \omega) 
\equiv \lambda f. (\lambda x.f (xx))(\lambda x.f (xx))}
\begin{enumerate}
    \vItem \iMbox{\mathrm{Y} \ f =_{\beta} f \ (\mathrm{Y} \ f)} i.e. 
    both \ul{$\beta$-reduce} to the same common term
    \vItem However \iMbox{f  \ (\mathrm{Y} \ f)} may not \textit{in general}
    \ul{$\beta$-reduce} to \iMbox{\mathrm{Y} \ f}
\end{enumerate}

\kwub{$Z$-combinator} => \iMbox{\mathrm{Z} \equiv \lambda f. 
(\lambda x.f (\lambda y. xxy))(\lambda x.f (\lambda y.xxy))},
\begin{enumerate}
    \vItem \iMbox{\mathrm{Z} \ f \ x = f \ (\mathrm{Z} \ f) \ x},
    i.e. next argument is defined \ul{explicitly}
    \vItem \stu{$\eta$-expansion} of \kwu{$Y$-combinator}, which will work for
    \stu{strict languages} (\kwu{$Y$-combinator} will stack overflow)
\end{enumerate}

\stub{Turing (\iMbox{\mathrm{U}}) and Theta (\iMbox{\mathrm{\Theta}}) combinators}:
\begin{enumerate}
    \vItem \iMbox{\mathrm{U} \equiv \lambda xy.y(xxy)} \&
    \iMbox{\Theta \equiv \mathrm{U}\mathrm{U}}

    \vItem These mirror the definitions of \ul{Mockingbird} and \ul{Omega} combinators

    \vItem \iMbox{\mathrm{\Theta} \ f \rightarrow_{\beta} f \ (\mathrm{\Theta} \ f)}, \\
    i.e. (unlike \stu{$Y$-combinator}) it actually \ul{$\beta$-reduces} 
\end{enumerate}

\hSep
